\documentclass{article}
\usepackage{amsmath,hyperref}
\title{Supply recovery with an old effect still live: two typed source fixtures}
\date{}
\begin{document}
\maketitle

\paragraph{Decision and limits.}
This I1 development block tests a small civil architecture with the same receiver admission rule on two different mission laws. The distinction is one official holder of a real report versus an additive power requirement for an already active cooling load. Independent physical mission data, platform faults and institutional authorization are not inferred from these synthetic source interpreters. The logical actor at the receiver can observe its instruments and current grants, not the world ID, unreceived acknowledgments or future fault. All writes and dispatches occur in simulation. This study does not freeze a confirmation family, spend the independent mine case, or label the inventory covered.

\paragraph{A shared receiver with two real model consequences.}
In one mission the old recorder P must remain the sole official custodian until an actually seated dock has completed a two-tick transfer. Receipt of a command does not change the holder. A local instrument measures usable docks; a shared recording head can prevent both from taking custody even when their control inputs accept a request. This source supports a mandatory report read at tick 6, so settlement at that tick is still useful. In the other mission a temporary pack supplies two units of modeled cooling continuously until its physical expiry. The alternate pickup path requires one available base-bus unit plus one available unit from its selected auxiliary. The two auxiliary paths feed the same exclusive selector; their availability cannot be added. Successful startup establishes actual continuous draw. Because one physically supplied replacement tick is owed before 6, the latest useful completion is 5. It is not enough to install a pending substitute at 6. When the old pack is about to die, an authorized protected shutdown loses cooling service but remains a separate safe contingency in this declared fixture.

Primary loss occurs at tick 0; during the positive envelope there is no usable RF report. A local setup operation occupies a diagnostic slot, and an onsite source report can arrive at 3 after observing the physical second loss at 2. A usable report supplies current route-specific readiness and a guaranteed interval; a limit crossing at 4 truncates that certificate even if the instantaneous report still looks healthy at 3. All source guarantees in this branch are limited to fixed physical law and finite time. The shared receiving rule accepts an authorized effect start only when the preparation interval fits within actual readiness, validity and continuing service of the old supplier, and finishes by the mission's latest useful effect cut. A source-specific instrument interpreter supplies the strength of evidence; it does not supply the joint action or a readiness oracle. Actual custody and live power are checked by separate physical event replays. The expiry that supports old work and the last useful effect cut must come from the actual mission or be conservatively certified; this pilot declares these inputs.

\paragraph{Fault-order, restart and mask controls.}
The main positive family uses cutoff 5 for the old supplier, one local diagnosis, second loss at 2 leaving either lawful standby, and two uninterrupted ticks to transfer from 3. There is an explicit policy that preserves the older duty while performing new mandatory useful service for every admitted single-standby loss, in both source laws. A report arriving a tick later cannot satisfy that same floor. If a common cause disables the recording head or the base bus, the policy can settle the older duty in a protected physical stop, but does not complete the mission. Two auxiliary controllers in the power example can each be correct that their own one-unit source is up; their claims do not imply the required two-unit output of a legally selected pickup. The checked joint premise is the pair of simultaneously consumed supplies with one permitted selection. A delayed second fault at tick 4 during acquisition yields a distinct boundary: both appear available at 3, but the short certificate cannot authorize a transfer finishing at 5. Extending the status-only interval by hand produces a physical failure. This ablation is not a fair performance comparison.

The restart pairs retain the same idle status with and without an older live claim. The independent optional work is allowed only on a scope-correct physical inspection of the old recorder seat or the discharge state, not from the shared idle screen. The optional control is charged and its actual event is replayed on the same histories as the recovery. Incorrectly treating masked status as clearance makes it start on an occupied slot. The positive mandatory recovery policies in this fixture need not complete that optional work.

\paragraph{Positive late-fault envelope with charged simultaneous preparation.}
A separate family has the second loss genuinely overlapping preparation: both alternatives have to begin their first setup step at 3 and actual loss cuts one at 4. Because the final physical effect requires a second setup step by 5, preparing only one candidate fails to guarantee one successful replacement in both possible loss histories. Preparing both with two real staging slots, then selecting exactly one accepted, locally verified survivor at 4, restores the service floor. The late local report belongs to the actuator-side selector and is not available to the earlier dispatcher. Without that report or authorization, no positive service is credited even if a viable physical configuration exists. The earlier operation has not released or finished during the two tentative setups.

For custody, pre-staged write fragments are not official custody or a valid delivered report. A single physical commitment at 5 moves official possession once and the unused candidate is discarded. This adapter is not admissible in a mission prohibiting those preparatory copies. For cooling, the staged setups consume independent low-power switching slots but cannot both energize the exclusive load pickup; actual selection draws one auxiliary and its physically available base. Therefore transfer changes the concrete type and eventual predicate while the receiving policy (prepared, authorized and presently usable, then one commit) and its proof of availability remain unchanged. Protection under a common cause refuses or performs safe shutdown; it cannot count as the positive service floor. A different fault cutting both physically live candidates invalidates the one-loss premise. A subsystem claiming more autonomy needs stronger source evidence and operational testing.

\paragraph{Evidence and remaining gates.}
Direct scenario-response enumeration shares the same action set, timings and local reports. It ties both receiver implementations. Parallel staging improves the single-staging reference of its explicitly restricted size and uses one extra real slot; a stronger plan that is allowed the same redundant preparation and local selector also obtains the guarantee. We therefore do not claim an advantage over the strongest established architecture. \href{https://arxiv.org/html/2103.11470}{NeBula} supplies a real baseline with mission executive and rescheduling; the previously run OPTIC and CSTN slices validate some admitted timing reductions, not the full policy or instrument implementation here. \href{https://openrepository.aut.ac.nz/items/9104ed9b-63a5-4975-8063-9034ed8efa12/full}{SPA-CBPA} supplies proactive successor selection with fallback reassignment; its own experimental completion convention and possession/supply interface must be translated before a matched benchmark. Additional source-specific validation and a frozen, independently selected changed-dependency mission are still needed. No platform certification, complete I1 inventory coverage, or I2 result is asserted.

The replay includes 768 main paired reference tests, 16 old-right timelines and 216 late staged tests, plus the explicit scope, shared-supply and late-certificate refusals. Four Lean lemmas check interval containment and the two-candidate bound. They do not formalize the source instruments. It is a finite proof that the local event policies suffice in the stated envelope, not a general distributed scheduling theorem.
\end{document}
